\chapter{Symbolic Expressions and Markers}
\label{ch:symbols}

\newthought{Many places in a task card} take a small math \emph{expression} rather than a plain
number---likelihood terms, derived quantities, selection conditions, and input mappings. This
chapter lists the symbols you can use in those expressions, and the path markers that keep your
file paths portable.

\section{Where expressions are used}
\label{sec:expr-where}

You will write expressions in:

\begin{itemize}
  \item \yamlkey{Sampling.LogLikelihood} (Chapter~\ref{ch:loglikelihood});
  \item \yamlkey{Sampling.selection} (the accept/reject condition);
  \item calculator input actions (Chapter~\ref{ch:calculators});
  \item \yamlkey{Operas} module inputs (Chapter~\ref{ch:operas}).
\end{itemize}

\section{Built-in constants}
\label{sec:expr-constants}

Three constants are always available:

\begin{description}[style=nextline, leftmargin=1.2em]
  \item[\code{Pi}] the number $\pi$.
  \item[\code{E}] Euler's number $e$.
  \item[\code{Inf}] infinity.
\end{description}

\noindent For example, the expression \expr{x * Pi} multiplies the variable \yamlkey{x} by $\pi$.

\section{Built-in functions}
\label{sec:expr-functions}

The most useful built-ins are the likelihood helpers and the common math functions
(Table~\ref{tab:builtins}).

\begin{table}[ht]
  \footnotesize
  \begin{tabular}{@{}lp{0.56\linewidth}@{}}
    \toprule
    \textbf{Function} & \textbf{Meaning} \\
    \midrule
    \expr{LogGauss(x, mean, err)} & log of a Gaussian --- the usual likelihood term \\
    \expr{Gauss(x, mean, err)}, \expr{Normal(x, mean, err)} & the Gaussian itself \\
    \code{sqrt}, \code{log}, \code{exp} & square root, natural log, exponential \\
    \code{sin}, \code{cos}, \code{tan} & trigonometric functions \\
    \code{Abs}, \code{Min}, \code{Max} & absolute value, minimum, maximum \\
    \bottomrule
  \end{tabular}
  \caption{Common built-in functions.}
  \label{tab:builtins}
\end{table}

\noindent Hyperbolic and inverse-trigonometric functions from \code{sympy} are available too. A
typical likelihood term reads \expr{LogGauss(mh, 125.09, 3.0)}.

\section{Your own and Operas functions}
\label{sec:expr-user}

Two more sources of functions become available when you configure them:

\begin{description}[style=nextline, leftmargin=1.2em]
  \item[\yamlkey{Utils} functions] Interpolation tables and helpers you declare yourself
        (Chapter~\ref{ch:utils}). For instance a tabulated limit curve \expr{XenonSD2019(mN1)}.
  \item[\yamlkey{Operas} functions] Operators exposed through \yamlkey{Operas.Functions}, when
        \Operas\ is installed (Chapter~\ref{ch:operas}).
\end{description}

\section{What an expression can read}
\label{sec:expr-symbols}

Inside any expression you may use:

\begin{itemize}
  \item the names of your scan variables (e.g. \code{m0}, \code{tanb});
  \item observables produced by earlier workflow modules (e.g. \code{mh}, \code{Omega\_h2});
  \item the built-in constants and functions above;
  \item your \yamlkey{Utils} functions and any whitelisted \yamlkey{Operas} functions.
\end{itemize}

\begin{jwarn}
A name is only usable \emph{after} it exists. If an expression reads an observable, the module
that produces it must run earlier in the workflow.
\end{jwarn}

\section{Selection conditions}
\label{sec:expr-selection}

A \yamlkey{selection} is a true/false expression used to accept or reject a point before the
heavy work runs:

\begin{yamlwin}
Sampling:
  selection: "(M1 > 0) & (TanBETA < 50)"
\end{yamlwin}

\noindent Use \code{\&} for ``and'' and \code{|} for ``or''. The expression must evaluate to a
strict boolean (true or false).

\section{Path markers}
\label{sec:expr-markers}

Path markers are not math, but they belong to the same authoring vocabulary. The one you will use
constantly is \marker{\&J}.

\begin{description}[style=nextline, leftmargin=1.2em]
  \item[\marker{\&J}] the project (task) root. If your card is \file{PROJECT/bin/task.yaml}, then
        \marker{\&J} resolves to \file{PROJECT/}. Use it for every project-local path, e.g.
        \code{save\_dir: "\&J/outputs"} or \code{source: "\&J/deps/program"}.
\end{description}

\noindent A leading \code{\textasciitilde} is expanded to your home directory. Plain relative
paths resolve from the project root, while \code{./} and \code{../} resolve relative to the
\textsc{yaml} file's own directory.

\begin{jnote}
\marker{\&J} is the only path marker the runtime resolves. For anything a scan would otherwise
read from outside the project, keep a copy under the project (for example \marker{\&J/deps/...})
and reference that---exactly what the built-in quick-starts do.
\end{jnote}

\section{Placeholders inside commands}
\label{sec:expr-placeholders}

Calculator and library commands accept a second kind of substitution: \emph{placeholders} written
with \code{\$\{...\}}, which expand to config values in scope.

\begin{description}[style=nextline, leftmargin=1.2em]
  \item[\code{\$\{path\}}, \code{\$\{source\}}] the calculator's runtime path and source path.
  \item[\code{\$\{...\}}] any other configuration value in scope.
  \item[\token{@PackID}] a per-instance id, giving each parallel sample its own working directory
        (Chapter~\ref{ch:core-concepts}).
\end{description}

\noindent For example: \code{installation: ["cp -r \$\{source\}/* \$\{path\}"]}.

\section{Practical advice}
\label{sec:expr-advice}

\begin{itemize}
  \item Start simple and build one observable at a time.
  \item Prefer clear names over clever shorthand.
  \item If an expression depends on a module output, make sure that module runs first.
\end{itemize}
